\documentclass{situationist-preprint}

\title{The Open Question Argument Revisited}
\msid{MTE-01-b}
\msfamily{MTE-01}
\msdiscipline{Metaethics}
\mstype{Companion manuscript}
\mskeywords{open question argument; normativity; triviality objection; a posteriori identity; moral naturalism}

\begin{document}
\maketitle

\begin{abstract}
MTE-01-a argues that moral facts are natural facts about idealized reactions and that Moore's open question argument fails, because its bridge from openness to distinctness of reference is refuted by a posteriori identities. Defenders of the argument have since recast it: the openness of moral questions is said to reflect the normativity of moral concepts, which no natural property can share. We examine the two strongest versions of this reply, the argument from normativity and Parfit's triviality objection. We show that the first restricts Moore's bridge without making it true, and that the second, if sound, would equally refute the identity of water and H$_2$O. The openness that the revised arguments rely on is predicted by the account of MTE-01-a.
\end{abstract}

\section{From synonymy to normativity}\label{s:norm}

Moore's original argument \citep[\S13]{Moore1903} concluded from the openness of every question of the form \enquote{Is a thing that is $N$ good?} that goodness is not the property expressed by any natural concept $N$. MTE-01-a shows that the argument is invalid as it stands (Proposition~5.2 of MTE-01-a) and that the premise needed to make it valid, Moore's bridge (Principle P2 of MTE-01-a), is false, since the question whether water is H$_2$O is open and the two concepts refer to one kind (Assumption~5.3 and Proposition~5.4 of MTE-01-a). This reply goes back to \citet{Frankena1939} and to the paradox of analysis \citep{Langford1942}, and it is now widely accepted against the argument in its original form.

The argument has nonetheless been thought to survive. \citet{DGR1992} suggest that the persistent openness of evaluative questions is best explained by the normativity of the concepts involved: to judge that something is good is to take oneself to have a reason to favor it, and no amount of natural information closes that question. On this view, the water case is beside the point. Natural kind concepts are not normative, and the bridge needs only to hold where one of the concepts is.

\begin{principle}[Normative bridge]\label{p:nb}
If the question whether a thing that is $C$ is $D$ is open, $D$ is a normative concept and $C$ is not, then $C$ and $D$ refer to different properties.
\end{principle}

Principle~\ref{p:nb} is not refuted by Assumption~5.3 of MTE-01-a, and the revised argument from it to non-naturalism is valid. The question is whether it is true.

\section{Normativity as a role clause}\label{s:role}

On the account of MTE-01-a, the reference of \enquote{wrong} is fixed by a role (Definition~3.4 and Principle P1 of MTE-01-a). One clause of the role is that the property bears on what rational agents do: an agent who takes an act to have it takes herself to have decisive reason against the act, and, if she is rational, she is moved accordingly, as MTE-07 argues. The normativity of the concept is this clause. It is a feature of the concept, which the role articulates, and not an additional feature that the property must have over and above playing the role.

\begin{proposition}\label{p:open}
On the account of MTE-01-a, for every natural concept $N$, including the concept of idealized disapproval, the question whether an act that is $N$ is wrong is open; and this is consistent with the identification of wrongness with idealized disapproval (Theorem~4.2 of MTE-01-a).
\end{proposition}

\begin{proof}
A speaker who possesses \enquote{wrong} possesses the role, and a speaker who possesses $N$ possesses a concept of a natural property. Whether the property expressed by $N$ plays the role depends on whether it satisfies the clauses of the role and is its unique realizer; for idealized disapproval this is Lemma~4.1 of MTE-01-a, which rests on the Uniqueness, Naturalness, Closure and Convergence assumptions (Assumptions~3.6--3.9 of MTE-01-a). These are substantive claims about human reactions and their convergence, not truths that reflection on the two concepts reveals. So the question is open. Openness concerns what reflection on concepts settles; Theorem~4.2 concerns which property the concepts pick out. The interpretation used in the proof of Proposition~5.2 of MTE-01-a, with \enquote{wrong} in place of \enquote{good}, satisfies both.
\end{proof}

Proposition~\ref{p:open} shows that Principle~\ref{p:nb} is not supported by the openness of moral questions, since the account of MTE-01-a predicts the openness while denying the distinctness of reference. The defender of the bridge must hold, further, that a normative concept cannot pick out its referent by a role whose realizer is found a posteriori. That is not something the openness of questions can show; it is the conclusion of the argument restated as a premise.

\section{The triviality objection}\label{s:triv}

\citet[vol.~2]{Parfit2011} presses a different point. If being wrong were the same property as being ideally disapproved, then the claim that wrong acts are ideally disapproved would state that ideally disapproved acts are ideally disapproved. But the claim is informative and is meant to guide action; on the naturalist view it would be trivial, and the normative content of moral claims would be lost.

The objection proves too much. If it were sound for normative concepts, a parallel objection would be sound for any a posteriori identity: if water is H$_2$O, then the claim that water is H$_2$O states that H$_2$O is H$_2$O. The reply in the water case is that the claim expresses an identity under two modes of presentation, and that its informativeness lies in the modes of presentation, not in the property \citep{Kripke1980}. The same reply is available here. That wrong acts are ideally disapproved is informative because \enquote{wrong} presents the property through its role, including its bearing on what to do, and \enquote{ideally disapproved} presents it through the reactions that realize the role. Nothing normative is lost: the property presented is the property that gives rational agents decisive reason against the act, which is just what the normative concept says of it. MTE-11 holds that reasons are primitive and that the other normative notions are defined in terms of them; that is a claim about the order of definition among normative concepts, and it is consistent with the a posteriori identity of what those concepts pick out with natural facts.

\section{Objections and replies}\label{s:obj}

\begin{objection}
Normativity is not a role. It is a sui generis authority that facts about reactions, however idealized, cannot have \citep{Enoch2011}.
\end{objection}

\begin{reply}
The objection asserts what the open question argument was meant to establish. If it is offered as a datum, the datum available is that competent speakers cannot settle by reflection whether any natural property has this authority; and that, by Proposition~\ref{p:open}, the account of MTE-01-a predicts. An argument from openness cannot then add anything to the bare assertion.
\end{reply}

\begin{objection}
Even an ideal observer could intelligibly ask whether what she would disapprove is wrong. The question remains open at the limit of inquiry, which a posteriori identities do not do.
\end{objection}

\begin{reply}
The ideal observer of MTE-01-a is ideal in information, imagination, reasoning and impartiality, not in her knowledge of metaethics. She may ask the question, as a chemist ignorant of the history of chemistry may ask whether water is H$_2$O. What closes the question is not a further idealization of the observer but the evidence that idealized disapproval satisfies the role, which is the argument of MTE-01-a, Section~4.
\end{reply}

\section{Conclusion}

The open question argument in its revised forms trades the false principle that openness excludes coreference for the principle that it excludes coreference where one concept is normative. The restricted principle has no independent support: on the role account, the normativity of moral concepts produces exactly the openness that the argument observes, while the concepts refer to a natural property. The triviality objection fails for the same reason that the parallel objection to the identity of water and H$_2$O fails. Moore's argument, revisited, does not tell against the naturalist realism of MTE-01-a.

\begin{thebibliography}{99}\small
\bibitem[Darwall et~al.(1992)]{DGR1992} Darwall, S., Gibbard, A. and Railton, P. (1992). Toward fin de si\`ecle ethics: Some trends. \emph{The Philosophical Review}, 101(1), 115--189.
\bibitem[Enoch(2011)]{Enoch2011} Enoch, D. (2011). \emph{Taking Morality Seriously: A Defense of Robust Realism}. Oxford University Press.
\bibitem[Frankena(1939)]{Frankena1939} Frankena, W.~K. (1939). The naturalistic fallacy. \emph{Mind}, 48(192), 464--477.
\bibitem[Kripke(1980)]{Kripke1980} Kripke, S.~A. (1980). \emph{Naming and Necessity}. Harvard University Press.
\bibitem[Langford(1942)]{Langford1942} Langford, C.~H. (1942). The notion of analysis in Moore's philosophy. In P.~A. Schilpp (ed.), \emph{The Philosophy of G.~E. Moore}. Northwestern University Press.
\bibitem[Moore(1903)]{Moore1903} Moore, G.~E. (1903). \emph{Principia Ethica}. Cambridge University Press.
\bibitem[Parfit(2011)]{Parfit2011} Parfit, D. (2011). \emph{On What Matters}, vols.~1--2. Oxford University Press.
\end{thebibliography}

\end{document}
